\documentclass[twocolumn]{aastex631}
\begin{document}
\title{The reionization ionizing-photon-budget shortfall for star-forming galaxies is not robust to systematics at z\~{}7 (jwsthigh systematic corner)}
\author{NebulaMind Lab (autonomous pipeline)}
\affiliation{NebulaMind Lab, \url{https://lab.nebulamind.net}}
\begin{abstract}
Reionization ionizing-photon-budget reconciliation at z\~{}7: star-forming galaxies require f\_esc=0.048 (+0.046/-0.023) to close the budget (Madau-Dickinson SFRD, log xi\_ion=25.5+/-0.15, clumping C=2-5, JWST-SFRD tail), versus indirect-proxy-inferred f\_esc=0.062 (+0.110/-0.039) from LzLCS O32/beta calibrations. Median delta(required-inferred)=-0.012 dex-frac (16-84\%: -0.120 to +0.042); 42\% of the systematic MC shows a shortfall. Conclusion: the budget CLOSES within the systematic, and the sign holds under both O32 and beta calibrations. The result is bounded by the xi\_ion x clumping x proxy-calibration systematic, not by statistics. Generated autonomously from public data (jwst) via the NebulaMind Lab runner: a bounded, reproducible, descriptive study.
\end{abstract}
\section{Introduction}
Reconciling the ionizing photon budget during reionization has been a long-standing challenge in understanding the early universe. Recent studies have highlighted potential discrepancies between the observed star formation rate density (SFRD) and the required ionizing photons to sustain reionization [Muñoz2024, Davies2021]. This tension raises questions about the contribution of star-forming galaxies to the photon budget and the need for alternative sources or mechanisms. To address this issue, we revisit the ionizing photon budget using a literature-anchored approach.
\section{Data and method}
Our method relies on the Madau \& Dickinson (2014) analytic fitting function for cosmic SFRD, combined with published values of xi\_ion and O32/beta f\_esc proxy calibrations from LzLCS [Chisholm+22, Flury+22; Simmonds+24]. We adopt a systematic approach to calculate the ionizing photon budget at z\~{}7, considering the effects of clumping factor (C=2-5) and the tail of the SFRD distribution. This allows us to assess whether star-forming galaxies alone can account for the required ionizing photons during reionization.
\section{Result}
Our calculation shows that star-forming galaxies require an escape fraction f\_esc = 0.048 (+0.046/-0.023) to close the ionizing photon budget at z\~{}7, assuming the Madau-Dickinson SFRD and log xi\_ion=25.5±0.15. This value is compared to the indirect-proxy-inferred f\_esc = 0.062 (+0.110/-0.039) from LzLCS O32/beta calibrations. The median difference between required and inferred escape fractions is -0.012 dex-frac, with a range of -0.120 to +0.042 (16-84\% confidence interval). Notably, 42\% of our systematic Monte Carlo simulations indicate a shortfall in the photon budget.
\begin{figure}\centering\includegraphics[width=\columnwidth]{result.png}\caption{The reionization ionizing-photon-budget shortfall for star-forming galaxies is not robust to systematics at z\~{}7 (jwsthigh systematic corner)}\end{figure}
\section{Caveats}
It is essential to acknowledge the limitations of this study. Our approach relies on automated, single-selection, and uncalibrated measurements, which may introduce biases and uncertainties. The adopted xi\_ion value and O32/beta calibrations are subject to systematic errors, potentially affecting our results. Additionally, the clumping factor and SFRD tail assumptions can impact the photon budget calculation. Therefore, while our study provides a reconciliation of the ionizing photon budget within the systematic uncertainties, further observations and refined calibrations are necessary to confirm these findings.
\section*{References}
\footnotesize
[Muoz2024] Muñoz et al. (2024). Reionization after JWST: a photon budget crisis?. bibcode 2024MNRAS.535L..37M \par [Park2022] Park et al. (2022). Calibrating excursion set reionization models to approximately conserve ionizing photons. bibcode 2022MNRAS.517..192P \par [Duncan2015] Duncan et al. (2015). Powering reionization: assessing the galaxy ionizing photon budget at z < 10. bibcode 2015MNRAS.451.2030D \par [Davies2021] Davies et al. (2021). The Predicament of Absorption-dominated Reionization: Increased Demands on Ionizing Sources. bibcode 2021ApJ...918L..35D \par [Madau2017] Madau (2017). Cosmic Reionization after Planck and before JWST: An Analytic Approach. bibcode 2017ApJ...851...50M

\end{document}
